% Research statement. Compile with cv/compile.ps1 (same chrome as the CV).
\documentclass[10pt,letterpaper]{article}
\usepackage{hyperref}
\usepackage{geometry}
\usepackage[T1]{fontenc}
\usepackage[sc,osf]{mathpazo}
\def\name{Regi Kusumaatmadja}
\def\footerlink{https://regikusumaatmadja.com/}
\hypersetup{
  colorlinks = true,
  urlcolor = black,
  pdfauthor = {\name},
  pdfkeywords = {economics, industrial organization, innovation, applied econometrics},
  pdftitle = {\name: Research Statement},
  pdfsubject = {Research Statement},
  pdfpagemode = UseNone
}
\geometry{
  body={6.5in, 8.5in},
  left=1.0in,
  top=1.25in
}
\pagestyle{myheadings}
\markboth{}{Research Statement}
\thispagestyle{empty}
\usepackage{sectsty}
\sectionfont{\rmfamily\mdseries\Large}
\subsectionfont{\rmfamily\mdseries\itshape\large}
\setlength\parindent{0em}
\setlength\parskip{0.55em}
\begin{document}
\centerline{\huge Research Statement}
\vspace{0.2in}

I am an applied economist working in industrial organization and the economics of innovation. My research asks how and why firms allocate their research: the mode of that research, why they enter certain areas, and how government assistance and changes in market structure affect the scale and composition of their portfolios. I combine applied theory with applied econometrics, using production and innovation surveys, patents, award records, and merger files.

\textbf{Internal and External R\&D: An analysis of costs and benefits.}
I study in-house research, contracted-out research, both, and neither. Using Dutch production and innovation surveys between 2000 and 2020, and focusing on the ICT industry, I document an increasing share of firms that perform R\&D and the four-way mix of those modes. To rationalize the mix, I build and estimate a dynamic discrete-choice model of R\&D with mode-specific investment costs. The cost of contracted-out R\&D is much higher than the cost of in-house R\&D, which explains the small share of firms that contract out alone. Doing both is cheaper than the sum of the two costs: in-house R\&D offsets almost all of the extra cost of going outside. I read that extra cost as a hold-up problem with the partner. The offset is absorptive capacity: a firm that already does in-house R\&D keeps more of the collaboration, so the two activities cost less together than apart. Productivity gains are similar across modes, so the mix of activities is driven by these costs. A uniform subsidy that does not favor a mode raises the share of R\&D-active firms and firm value by more than a subsidy aimed only at in-house R\&D.

\textbf{Entry into Innovation Areas.}
I then ask where firms put that research. The paper studies multi-area entry by publicly listed firms in the United States electronics industry from 1990 to 2019. I apply Latent Dirichlet Allocation to over 630,000 patent titles and abstracts to endogenously define twenty innovation areas, and I merge these with market-implied patent valuations. Three patterns stand out. Innovation occurs across multiple domains at once. The number of active competitors varies across areas. The industry focus shifted from physical hardware to connectivity technologies. To explain these patterns, I estimate a static simultaneous-move game of firm entry using moment inequalities, which accommodates multiple equilibria and the combinatorial choice set. The 99 percent confidence sets support four conclusions. Firms enter areas where expected patent valuations are higher. Holding those valuations fixed, additional rivals reduce expected profit. Firms also benefit from occupying more areas at once. A large per-area entry cost, together with rival congestion, still limits how many areas a firm enters. Starting from observed portfolios, between 13 and 16 percent of firm-years already have at least one extra area that looks worth occupying at the estimated costs. A 25 percent reduction in the per-area cost raises that share to between 22 and 27 percent, and a 50 percent reduction raises it to between 33 and 41 percent. A 50 percent reduction in how costly rivals are raises that share to between 15 and 21 percent. In a case study of targeted cost reductions, entry rises in semiconductor areas such as Semiconductor Memory Architecture and Semiconductor Lithography and Fabrication Processes. The offset is lower entry in other fields, including Wireless Base-Station and Mobile Communications.

\textbf{The Impact of Government Financial Assistance on the Scope and Direction of Firm Innovation}, with Costanza Cincotta.
Governments spend hundreds of billions of dollars each year to support private-sector innovation, but it is unclear whether this money expands existing research, redirects it, or reallocates it away from other firms. We develop a two-stage model in which firms first choose an R\&D portfolio and then compete in prices \`a la Salop. Financially constrained firms underinvest, and even unconstrained firms under-provide quality because product-market competition leaves a gap between private and social returns. Liquidity assistance can relax the cash constraint, while the quality gap remains. Cheaper cash can expand how much a recipient patents and how many technology areas it covers. If the firm simply scales up the research it already conducts, the composition of that research does not change. Assistance that is targeted toward particular activities can redirect the portfolio. Nearby rivals may expand through knowledge pooling or contract through business stealing. We take these implications to a panel of U.S.\ public firms from 2000 to 2021, combining Good Jobs First awards with patent data. Scope measures how much a firm patents and how many technology areas it covers. Direction measures whether those patents align with frontier or peer technologies. After reweighting untreated firms so that they match recipients on pre-treatment characteristics, we estimate a stacked difference-in-differences. Distant rivals provide the preferred reading of the treated firm's own response. Relative to distant industry controls, intensive scope rises by about 10.5 percent after an award. Relative to firms with little product-market overlap, the same contrast is about 3 percent. Direction shifts toward the industry's most-cited technologies, and the shift is strongest for cash grants. When nearby means high text-based product-market overlap, those nearby rivals' patenting scale does not move. When nearby means the same two-digit industry, industry-code neighbors expand as well, the knowledge-pooling case.

\textbf{Mergers and Innovation: An Exploration of Scope and Direction of Innovation}, with Sabien Dobbelaere and Jos\'e Luis Moraga-Gonz\'alez.
How do mergers change the scale and composition of firm innovation? We study U.S.\ public mergers from 1980 to 2020 using combined acquirer-target patent portfolios. We measure \emph{scope}, the scale of the joint portfolio, and \emph{direction}, its alignment with highly cited patents. A compact portfolio model isolates three forces: product-market business stealing, technological knowledge pooling, and fixed costs of activating innovation lines. The model leaves the average change in the split between shared and unique lines unsigned, and it leaves the average change in direction unsigned as well. Empirically, we compare treated merger pairs to matched control pairs. On the matched sample, joint-portfolio scope is larger after mergers than for matched controls. The difference is substantially larger when pre-merger technologies are more related, especially among non-horizontal deals. We find no strong evidence of reallocation: innovation activity expands after mergers, without a clear rebalancing between technology classes shared by both firms and classes unique to one. Joint portfolios also show higher alignment with highly cited patents relative to matched controls. Acquirer-only R\&D measures overstate post-merger input growth relative to joint-entity measurement.

These papers treat the allocation of research as a set of discrete choices over modes, areas, and portfolios. Costs, rivals, cash, and ownership then move both how much firms invent and what they invent.

\vspace{0.15in}
\begin{center}
  \begin{footnotesize}
    Last updated: \today
  \end{footnotesize}
\end{center}
\end{document}
